% Source: openai/math@adc7f124, 06-transport.tex, lines 257--286;
% HistoryMeanTree.lean fixes the actual binary parenthesizations.
% Checked identities: historyMeanTree_eval_succ, fill_interpRoot_zero,
% fill_interpRoot_one, charge_interpRoot_zero, charge_interpRoot_one.
% Physical metric identifications: FillTransportEndpoints.lean at f84f568b,
% actualFillData_rootPath_one_eq_charge_zero and
% charge_rootPath_one_eq_actualFillData_zero. These repaired candidates remain
% kernel-pending; this explanatory component adds no proof-status markers.
% The notation h \frown c below is extendHistory h c in the source.

Write $T_k$ for the recursive history tree and $C$ for the fixed charge-choice
tree of Definition~\ref{def:al_scan_history_mean_tree}. Fix positive local
dimensions, a replica count, and $t\ge0$. Let $A_h$ be the completed-history metric and
$\widetilde A_h$ its pre-charge metric. Write $h\frown c$ for the history
obtained by appending charge $c$, and $T[A]$ for a tree evaluated at its
terminal metrics. The endpoint substitutions are
\begin{align}
    T_k[A_h] & \xrightarrow{\mathrm{fill}} T_k[\widetilde A_h]
    \xrightarrow{\mathrm{charge}} T_k\bigl[{C[A_{h\frown c}]}\bigr]
    =T_{k+1}[A_{h\frown c}].
    \label{eq:al_scan_recursive_substitution}
\end{align}

\begin{figure}[ht]
    \centering
    % Evaluation schematic, not a tensor contraction. Each box denotes a
    % whole fixed tree (T_k or C) or a terminal matrix (A_i or A_ij).
    % A dotted string abbreviates an evaluation path, directed toward the
    % root box. It is not an index: no contracted or open tensor legs occur,
    % so every panel has empty tensor boundary signature.
    % Two representative histories and two representative charges are drawn;
    % no claim is made that either actual finite tree has only two leaves.
    % A_i = A_(h_i); tilde A_i = tilde A_(h_i);
    % A_ij = A_(extendHistory h_i c_j). The two C boxes are the same fixed tree.
    % All labels are inside boxes. Empty grid rows and columns separate them;
    % no warning, label, bounding-box, or equation check is disabled.
    \tenkzkernel
    \begin{tenkzequation}
        \begin{tenkz}[rows={wire,wire,wire}, cols=3, bonds=none, west=none, east=none]
            \tn[at=(1,2), name=fillOldRoot, skin=box]{T_k}
            \tn[at=(3,1), name=fillOldZero, skin=box]{A_0}
            \tn[at=(3,3), name=fillOldOne, skin=box]{A_1}
            \tnwire[kind=string, stroke=dotted, dir=to, name=fillOldPathZero]{fillOldZero}{fillOldRoot}
            \tnwire[kind=string, stroke=dotted, dir=to, name=fillOldPathOne]{fillOldOne}{fillOldRoot}
        \end{tenkz}
        $\xrightarrow{\mathrm{fill}}$
        \begin{tenkz}[rows={wire,wire,wire}, cols=3, bonds=none, west=none, east=none]
            \tn[at=(1,2), name=fillNewRoot, skin=box]{T_k}
            \tn[at=(3,1), name=fillNewZero, skin=box]{\widetilde A_0}
            \tn[at=(3,3), name=fillNewOne, skin=box]{\widetilde A_1}
            \tnwire[kind=string, stroke=dotted, dir=to, name=fillNewPathZero]{fillNewZero}{fillNewRoot}
            \tnwire[kind=string, stroke=dotted, dir=to, name=fillNewPathOne]{fillNewOne}{fillNewRoot}
        \end{tenkz}
    \end{tenkzequation}
    \begin{tenkzequation}
        \begin{tenkz}[rows={wire,wire,wire}, cols=3, bonds=none, west=none, east=none]
            \tn[at=(1,2), name=chargeOldRoot, skin=box]{T_k}
            \tn[at=(3,1), name=chargeOldZero, skin=box]{\widetilde A_0}
            \tn[at=(3,3), name=chargeOldOne, skin=box]{\widetilde A_1}
            \tnwire[kind=string, stroke=dotted, dir=to, name=chargeOldPathZero]{chargeOldZero}{chargeOldRoot}
            \tnwire[kind=string, stroke=dotted, dir=to, name=chargeOldPathOne]{chargeOldOne}{chargeOldRoot}
        \end{tenkz}
        $\xrightarrow{\mathrm{charge}}$
        \begin{tenkz}[rows={wire,wire,wire,wire,wire}, cols=7, bonds=none, west=none, east=none]
            \tn[at=(1,4), name=chargeNewRoot, skin=box]{T_k}
            \tn[at=(3,2), name=chargeNewChoiceZero, skin=box]{C}
            \tn[at=(3,6), name=chargeNewChoiceOne, skin=box]{C}
            \tn[at=(5,1), name=chargeNewZeroZero, skin=box]{A_{00}}
            \tn[at=(5,3), name=chargeNewZeroOne, skin=box]{A_{01}}
            \tn[at=(5,5), name=chargeNewOneZero, skin=box]{A_{10}}
            \tn[at=(5,7), name=chargeNewOneOne, skin=box]{A_{11}}
            \tnwire[kind=string, stroke=dotted, dir=to, name=chargeNewHistoryZero]{chargeNewChoiceZero}{chargeNewRoot}
            \tnwire[kind=string, stroke=dotted, dir=to, name=chargeNewHistoryOne]{chargeNewChoiceOne}{chargeNewRoot}
            \tnwire[kind=string, stroke=dotted, dir=to, name=chargeNewPathZeroZero]{chargeNewZeroZero}{chargeNewChoiceZero}
            \tnwire[kind=string, stroke=dotted, dir=to, name=chargeNewPathZeroOne]{chargeNewZeroOne}{chargeNewChoiceZero}
            \tnwire[kind=string, stroke=dotted, dir=to, name=chargeNewPathOneZero]{chargeNewOneZero}{chargeNewChoiceOne}
            \tnwire[kind=string, stroke=dotted, dir=to, name=chargeNewPathOneOne]{chargeNewOneOne}{chargeNewChoiceOne}
        \end{tenkz}
    \end{tenkzequation}
    \caption{Schematic of the recursive scan trees in
        (\ref{eq:al_scan_recursive_substitution}). Here $A_i=A_{h_i}$,
        $\widetilde A_i=\widetilde A_{h_i}$, and
        $A_{ij}=A_{h_i\frown c_j}$. Dotted arrows abbreviate whole evaluation
        paths; only two histories and two choices per history are shown.
        The deterministic fill changes terminal metrics on the same tree.
        The charge attaches the same fixed choice tree at every history leaf
        and relabels $(h,c)$ by $h\frown c$, giving $T_{k+1}$.
        Every internal geometric mean retains its specified weight and
        parenthesization. No reassociation or replacement by an independent
        uniform average is used; the endpoint identities are those of
        Theorem~\ref{thm:al_scan_history_mean_tree_endpoints}.}
    \label{fig:al_scan_recursive_transport}
\end{figure}
